\documentclass[11pt]{article}

\usepackage{amsmath,amssymb}
\usepackage{xurl}
\usepackage[colorlinks=true,linkcolor=blue,citecolor=blue,urlcolor=blue]{hyperref}

% Shared commands for notes in docs/paper-gaps/.

\DeclareUnicodeCharacter{2080}{\ifmmode{}_{0}\else\textsubscript{0}\fi}
\DeclareUnicodeCharacter{2081}{\ifmmode{}_{1}\else\textsubscript{1}\fi}
\DeclareUnicodeCharacter{2082}{\ifmmode{}_{2}\else\textsubscript{2}\fi}
\DeclareUnicodeCharacter{2099}{\ifmmode{}_{n}\else\textsubscript{n}\fi}

\newcommand{\Lean}{\textsc{Lean}}
\ExplSyntaxOn
\NewDocumentCommand{\leanid}{m}{
  \ifmmode
    \textsf{\detokenize{#1}}
  \else
    \group_begin:
    \urlstyle{sf}
    \tl_set:Nn \l_tmpa_tl {#1}
    \tl_replace_all:Nnn \l_tmpa_tl {\_} {_}
    \exp_args:NV \path \l_tmpa_tl
    \group_end:
  \fi
}
\ExplSyntaxOff
\providecommand{\tr}{\operatorname{tr}}
\newcommand{\tnleanIssueBase}{https://github.com/LionSR/TNLean/issues/}
\newcommand{\tnleanPrBase}{https://github.com/LionSR/TNLean/pull/}
\newcommand{\tnleanDocBase}{https://sirui-lu.com/TNLean/docs/}
\newcommand{\ghissue}[1]{\href{\tnleanIssueBase#1}{GitHub issue \##1}}
\newcommand{\ghpr}[1]{\href{\tnleanPrBase#1}{PR \##1}}
\newcommand{\doclink}[1]{\href{\tnleanDocBase#1.html}{\path{#1}}}

% Dirac notation and the identity operator, as in blueprint/src/macros/common.tex.
% \providecommand keeps note-local definitions authoritative.
\usepackage{dsfont}
\providecommand{\ket}[1]{|#1\rangle}
\providecommand{\bra}[1]{\langle #1|}
\providecommand{\braket}[2]{\langle #1 | #2 \rangle}
\providecommand{\ketbra}[2]{|#1\rangle\!\langle #2|}
\providecommand{\Id}{\mathds{1}}
\providecommand{\one}{\mathds{1}}
\providecommand{\C}{\mathbb{C}}
\providecommand{\MN}[1]{M_{#1}(\C)}
\providecommand{\MD}{M_D(\C)}

% External referencing for paper sources.  The build helper writes sanitized
% auxiliary files under build/paper-aux/ and excludes duplicated source labels.
\usepackage{xr}

% Import a generated paper auxiliary file with a caller-supplied label prefix.
% The candidate paths support builds from docs/paper-gaps/, the repository root,
% and build/paper-aux/ itself.
\newcommand{\importpaperaux}[2]{%
  \IfFileExists{../../build/paper-aux/#2.aux}{%
    \externaldocument[#1]{../../build/paper-aux/#2}%
  }{%
    \IfFileExists{build/paper-aux/#2.aux}{%
      \externaldocument[#1]{build/paper-aux/#2}%
    }{%
      \IfFileExists{#2.aux}{%
        \externaldocument[#1]{#2}%
      }{}%
    }%
  }%
}

% General external reference macros
\newcommand{\paperref}[2]{\ref{#1-#2}}
\newcommand{\papereqref}[2]{(\ref{#1-#2})}

% CPSV16 source (1606.00608 / MPDO-22-12-17-2.tex) external document and shortcuts
\importpaperaux{cpsv16-}{1606.00608/MPDO-22-12-17-2-tnlean}
\newcommand{\cpsvref}[1]{\paperref{cpsv16}{#1}}
\newcommand{\cpsveqref}[1]{\papereqref{cpsv16}{#1}}

% Verdict marker: \gapnote{<kind>}{<status>}. Kind is the policy
% classification (clarification, local-correction, scope-restriction,
% unfaithful, false-source, open-gap); status is open, wip (elimination
% underway), resolved, or historical. Machine-read by the site index;
% prints a verdict line.
\newcommand{\gapnote}[2]{\par\noindent\textbf{Verdict:} #1 (#2).\par}


\title{Full-Matrix Scope for the GNVW Support-Algebra Formalization}
\author{}
\date{2026-08-24}

\begin{document}
\maketitle
\gapnote{scope-restriction}{open}

\section{The general support-algebra statement}

Let $\mathcal B_1$ and $\mathcal B_2$ be finite-dimensional $C^*$-algebras,
and let $\mathcal A\subseteq\mathcal B_1\otimes\mathcal B_2$ be a
$C^*$-subalgebra. There is a smallest $C^*$-subalgebra
$\mathcal C_1\subseteq\mathcal B_1$ such that
\begin{align}
    \mathcal A
    &\subseteq\mathcal C_1\otimes\mathcal B_2,
    \label{eq:gnvw12_support_algebra_full_matrix_scope_support_containment}\\
    \mathcal C_1
    &=\operatorname{Spp}(\mathcal A,\mathcal B_1).
    \label{eq:gnvw12_support_algebra_full_matrix_scope_support_definition}
\end{align}
Choose a basis $(e_\mu)_\mu$ of $\mathcal B_2$. Every $a\in\mathcal A$ has a
unique expansion
\begin{align}
    a
    &=\sum_\mu a_\mu\otimes e_\mu.
    \label{eq:gnvw12_support_algebra_full_matrix_scope_coefficient_expansion}
\end{align}
The support algebra
$\operatorname{Spp}(\mathcal A,\mathcal B_1)$ is generated by all
coefficients $a_\mu$ obtained from all $a\in\mathcal A$. Its commutant inside
$\mathcal B_1$ is
\begin{align}
    \operatorname{Spp}(\mathcal A,\mathcal B_1)'
    &=\left\{b\in\mathcal B_1:
      b\otimes\mathbf{1}\in\mathcal A'\right\}.
    \label{eq:gnvw12_support_algebra_full_matrix_scope_commutant}
\end{align}

These are the three assertions of Lemma~7 in Section~7.1 of
Ref.~\cite{Gross2012IndexTheory}, after applying the necessary star-closure
correction to the commutant clause. That correction is recorded separately in
\path{docs/paper-gaps/gnvw12_support_algebra_star_closure.tex}. In the arXiv
v2 source, the statement and leastness argument appear at lines 1181--1191,
the coefficient-generation argument appears at lines 1188--1191, and the
commutant assertion appears at lines 1192--1208. The published DOI
\href{https://doi.org/10.1007/s00220-012-1423-1}
{10.1007/s00220-012-1423-1} and the arXiv v2 source are
third-party-verifiable anchors. The source label for the lemma is
\path{lem:spp}.

Applying the construction to both factors gives
\begin{align}
    \mathcal A
    &\subseteq
      \operatorname{Spp}(\mathcal A,\mathcal B_1)
      \otimes
      \operatorname{Spp}(\mathcal A,\mathcal B_2),
    \label{eq:gnvw12_support_algebra_full_matrix_scope_two_support_containment}\\
    \operatorname{Spp}(\mathcal A,\mathcal B_1)
      \otimes
      \operatorname{Spp}(\mathcal A,\mathcal B_2)
    &\subseteq\mathcal B_1\otimes\mathcal B_2.
    \label{eq:gnvw12_support_algebra_full_matrix_scope_ambient_containment}
\end{align}
The first inclusion is GNVW equation~(22), which follows after Lemma~7 and is
not one of that lemma's three clauses
\cite{Gross2012IndexTheory}. It appears at lines 1211--1219 of the arXiv v2
source with label \path{a2Supp}. Schumacher and Werner give the preceding
coefficient-space and generated-support-algebra constructions in
Section~4.3 of Ref.~\cite{SchumacherWerner2004QCA}; see lines 1137--1169 and
the source labels \path{a2supp} and \path{a2Supp}.

\section{The full-matrix formal statement}

Let $I$ and $J$ be finite sets, and let
\begin{align}
    \mathcal A
    &\subseteq M_{I\times J}(\mathbb C),
    \label{eq:gnvw12_support_algebra_full_matrix_scope_matrix_input}\\
    M_{I\times J}(\mathbb C)
    &\cong M_I(\mathbb C)\otimes M_J(\mathbb C).
    \label{eq:gnvw12_support_algebra_full_matrix_scope_matrix_factorization}
\end{align}
Assume that $\mathcal A$ is a star-subalgebra. This assumption implements the
resolved local correction described above. The remaining open scope concerns
only the extension from full complex matrix factors to arbitrary
finite-dimensional $C^*$-algebra factors.

For $X\in\mathcal A$, define its left and right coefficient matrices by
\begin{align}
    X^{\mathrm L}_{rs}(i,j)
    &=X((i,r),(j,s)),
    \label{eq:gnvw12_support_algebra_full_matrix_scope_left_coefficient}\\
    X^{\mathrm R}_{ij}(r,s)
    &=X((i,r),(j,s)).
    \label{eq:gnvw12_support_algebra_full_matrix_scope_right_coefficient}
\end{align}
The formal support algebras are the generated star-algebras
\begin{align}
    \operatorname{Spp}_{\mathrm L}(\mathcal A)
    &=\operatorname{alg}^*
      \left\{X^{\mathrm L}_{rs}:
      X\in\mathcal A,\ r,s\in J\right\},
    \label{eq:gnvw12_support_algebra_full_matrix_scope_left_support}\\
    \operatorname{Spp}_{\mathrm R}(\mathcal A)
    &=\operatorname{alg}^*
      \left\{X^{\mathrm R}_{ij}:
      X\in\mathcal A,\ i,j\in I\right\}.
    \label{eq:gnvw12_support_algebra_full_matrix_scope_right_support}
\end{align}

For a star-subalgebra $\mathcal D\subseteq M_I(\mathbb C)$, left-support
leastness is equivalent to coefficient containment:
\begin{align}
    \operatorname{Spp}_{\mathrm L}(\mathcal A)
    \subseteq\mathcal D
    \quad&\Longleftrightarrow\quad
    X^{\mathrm L}_{rs}\in\mathcal D
    \quad
    \text{for all }X\in\mathcal A,\ r,s\in J.
    \label{eq:gnvw12_support_algebra_full_matrix_scope_coefficient_leastness}
\end{align}
If
\begin{align}
    I
    &=\{0,\ldots,p-1\},
    \label{eq:gnvw12_support_algebra_full_matrix_scope_left_finset}\\
    J
    &=\{0,\ldots,q-1\},
    \label{eq:gnvw12_support_algebra_full_matrix_scope_right_finset}
\end{align}
then coefficient leastness is equivalent to tensor-factor leastness:
\begin{align}
    \operatorname{Spp}_{\mathrm L}(\mathcal A)
    \subseteq\mathcal D
    \quad&\Longleftrightarrow\quad
    \mathcal A\subseteq\mathcal D\otimes M_q(\mathbb C).
    \label{eq:gnvw12_support_algebra_full_matrix_scope_tensor_leastness}
\end{align}
The commutant theorem states
\begin{align}
    B\otimes\mathbf{1}_J\in\mathcal A'
    \quad&\Longleftrightarrow\quad
    B\in\operatorname{Spp}_{\mathrm L}(\mathcal A)'.
    \label{eq:gnvw12_support_algebra_full_matrix_scope_matrix_commutant}
\end{align}
The right-factor statements are analogous.

For subspaces $S\subseteq M_I(\mathbb C)$ and
$T\subseteq M_J(\mathbb C)$, define
\begin{align}
    S\boxtimes T
    &:=\operatorname{span}_{\mathbb C}
      \left\{A\otimes B:
      A\in S,\ B\in T\right\}.
    \label{eq:gnvw12_support_algebra_full_matrix_scope_boxtimes}
\end{align}
The formal counterpart of GNVW equation~(22) of Ref.~\cite{Gross2012IndexTheory} is the literal containment
\begin{align}
    \mathcal A
    &\subseteq
      \operatorname{Spp}_{\mathrm L}(\mathcal A)
      \boxtimes
      \operatorname{Spp}_{\mathrm R}(\mathcal A).
    \label{eq:gnvw12_support_algebra_full_matrix_scope_matrix_two_factor}
\end{align}
This statement places every element of $\mathcal A$ in the span of simple
tensors whose factors lie in the two support algebras. It is stronger in form
than merely asserting both one-sided containments.

The formal definitions and theorems are
\leanid{Matrix.leftSupportAlgebra},
\leanid{Matrix.rightSupportAlgebra},
\leanid{Matrix.left_support_algebra_le_iff},
\leanid{Matrix.left_support_algebra_le_iff_le_kronecker_submodule},
\leanid{Matrix.le_support_kroneckerSubmodule},
\leanid{Matrix.left_kronecker_embed_mem_centralizer_iff}, and their
right-factor counterparts in
\doclink{TNLean/QCA/BipartiteSupportAlgebra}.

\section{Why the specialization suffices for the MPU application}

The matrix-product-unitary Appendix observes that every one-site algebra is
isomorphic to a full matrix algebra,
\begin{align}
    \mathcal A_z
    &\cong M_d(\mathbb C),
    \label{eq:gnvw12_support_algebra_full_matrix_scope_mpu_site_algebra}
\end{align}
before invoking the GNVW support algebra; see lines 2320--2334 of
Ref.~\cite{Cirac2017MPU}. Each two-site factor is therefore isomorphic to
\begin{align}
    \mathcal A_z\otimes\mathcal A_{z+1}
    &\cong M_{d^2}(\mathbb C).
    \label{eq:gnvw12_support_algebra_full_matrix_scope_mpu_two_site_algebra}
\end{align}
After choosing these matrix coordinates, the support algebra
$\mathcal R_{2x}$ is an instance of
(\ref{eq:gnvw12_support_algebra_full_matrix_scope_coefficient_leastness}) and
(\ref{eq:gnvw12_support_algebra_full_matrix_scope_tensor_leastness}).
Consequently, the full-matrix theorem suffices for the support-algebra step in
the spin-chain argument of Ref.~\cite{Cirac2017MPU}.

The star-closed-input correction and its verification for the GNVW and MPU
applications are recorded in
\path{docs/paper-gaps/gnvw12_support_algebra_star_closure.tex}. They are not
part of the open scope considered here.

The restriction to full matrix factors is an open formalization boundary in
TNLean. It is not a mathematical gap in Lemma~7 or equation~(22) of
Ref.~\cite{Gross2012IndexTheory}.

\section{Elimination plan}

Removing the full-matrix restriction requires the same construction for
arbitrary finite-dimensional $C^*$-algebras $\mathcal B_1$ and $\mathcal B_2$
and a $C^*$-subalgebra
$\mathcal A\subseteq\mathcal B_1\otimes\mathcal B_2$. Choose a basis
$(e_\mu)_\mu$ of $\mathcal B_2$ and its dual basis $(\omega_\mu)_\mu$. Extract
the coefficients by
\begin{align}
    a_\mu
    &=(\operatorname{id}\otimes\omega_\mu)(a).
    \label{eq:gnvw12_support_algebra_full_matrix_scope_dual_coefficient}
\end{align}
Define $\operatorname{Spp}(\mathcal A,\mathcal B_1)$ as the
$C^*$-subalgebra generated by these coefficients.

The dual-basis expansion
(\ref{eq:gnvw12_support_algebra_full_matrix_scope_coefficient_expansion})
proves the leastness property
(\ref{eq:gnvw12_support_algebra_full_matrix_scope_support_containment}).
Expanding the commutator gives
\begin{align}
    [b\otimes\mathbf{1},a]
    &=\sum_\mu[b,a_\mu]\otimes e_\mu.
    \label{eq:gnvw12_support_algebra_full_matrix_scope_commutator_expansion}
\end{align}
This identity proves
(\ref{eq:gnvw12_support_algebra_full_matrix_scope_commutant}). Repeating the
construction for the second factor proves
(\ref{eq:gnvw12_support_algebra_full_matrix_scope_two_support_containment}).
The existing full-matrix statements should then follow by taking matrix units
as the chosen basis and using the associated coefficient functionals. This is
the direct coefficient argument of Ref.~\cite{Gross2012IndexTheory}; it
requires neither a bicommutant theorem nor a density theorem.

Until this generic construction is formalized, a source-labelled entry may
carry a completion marker only if it states the full complex matrix-factor
restriction explicitly.

\bibliographystyle{alpha}
\bibliography{references}

\end{document}
